\documentclass[letterpaper]{article}
\PassOptionsToPackage{unicode=true}{hyperref} % options for packages loaded elsewhere
\PassOptionsToPackage{hyphens}{url}
\usepackage{graphicx} % Required for inserting images


\def\tightlist{}


\graphicspath{ {../images/} }
\usepackage[margin=1in]{geometry}
\usepackage{tikz}
\usepackage[dvipsnames]{xcolor}
\usetikzlibrary{calc}




\usepackage[utf8]{inputenc}
\usepackage[T1]{fontenc}
\usepackage{lmodern}
\usepackage[english]{babel}

\usepackage{longtable}
\usepackage{booktabs}

% Useful packages
\usepackage{hyperref}
\usepackage{amsmath, amssymb}
\usepackage{fancyhdr}
\usepackage{setspace}
\usepackage{tocloft}
\usepackage{titlesec}


\usepackage{unicode-math}



% Code highlighting (Pandoc uses listings or minted, but listings is safer without shell escape)
\usepackage{listings}
\lstset{
  basicstyle=\ttfamily\small,
  breaklines=true,
  frame=single,
  backgroundcolor=\color[gray]{0.95}
}



%Modifiable Commands; defaults set to none
\newcommand{\titleDOC}{}
\newcommand{\authorDOC}{Daniel Sabalakov}
\newcommand{\dateDOC}{\today}
\newcommand{\classDOC}{}
\newcommand{\emailDOC}{danielsabalakov@gmail.com}
% DEFAULT QUOTE IS BY ISSAC NEWTON. IF YOU DON'T WANT IT, DO: quote: @none
\newcommand{\quoteDOC}{``\textit{Truth is ever to be found in the simplicity, and not in the multiplicity and confusion of things}'' - Issac Newton}

% Header/Footer
\pagestyle{fancy}
\fancyhf{}
\rhead{\thepage}
\lhead{\leftmark}





% %Title Arguments
% \title{\TITLEOFDOCUMENT}
% \author{\AUTHOROFDOCUMENT}
% \date{\today}


%
% ANYTHING BELOW HERE CAN GET PROCEDURALLY GENERATED
%
% BEGINNING OF DOCUMENT
%
%
%
\begin{document}

%titlepage








\renewcommand{\classDOC}{Trigonometry}

\renewcommand{\authorDOC}{Daniel Sabalakov}

\renewcommand{\titleDOC}{Unit 14 Review}





\begin{titlepage}
  \titleDOC
  \classDOC
  \dateDOC
  \authorDOC
\end{titlepage}

\section{Topics covered in quiz}\label{topics-covered-in-quiz}

\begin{itemize}
\tightlist
\item
  Graphing using Unit Circle
\item
  Solving equations with general solutions / restrictive intervals
\item
  Modeling sinusoidal functions given Maximum and Minimum
\end{itemize}

\section{Unit Circle}\label{unit-circle}

\begin{figure}[htbp]
	 \centering
	 \includegraphics[width=0.9\textwidth]{C:/Users/score/Downloads/ucirc.png}
\end{figure}

\(\newline\)

\section{Questions}\label{questions}

\subsection{Topic 1: Solving Equations with General Solutions / Restrictive Intervals}\label{topic-1-solving-equations-with-general-solutions-restrictive-intervals}

Solve the following equations. Assume the unknown variable is between {[}0,\(\frac{\pi}{2}\)) (In radians)

\begin{enumerate}
\def\labelenumi{\arabic{enumi}.}
\tightlist
\item
  \(2cos(2x)-\sqrt{3} = 0\)
\item
  \(4sin^2(\theta)-4=1\)
\item
  \(3tan(x)+\sqrt{3}=0\)
\item
  \(sec(u)\sqrt{3}=2\)
\item
  \textbf{Challenge:} \(2sin^2(z)*cot^2(z)=1\)
\end{enumerate}

Solve the following equations for their general solutions:

\begin{enumerate}
\def\labelenumi{\arabic{enumi}.}
\tightlist
\item
  \(\sin(x) = \frac{1}{2}\)
\item
  \(\cos(x) = -\frac{\sqrt{2}}{2}\)
\item
  \(\tan(x) = \sqrt{3}\)
\item
  \(\sin(x) = -\frac{\sqrt{3}}{2}\)
\item
  \textbf{Challenge:} \(2\cos^2(x) - 3\sin(x) = 0\)
\end{enumerate}

\subsection{Topic 2: Modelling sinusoidal functions given Maximum and Minimums}\label{topic-2-modelling-sinusoidal-functions-given-maximum-and-minimums}

Model the following three problems using Sine. Assume any point is either a maximum or a minimum point.

\begin{enumerate}
\def\labelenumi{\arabic{enumi}.}
\tightlist
\item
  \((0,0),(1,\pi)\)
\item
  \((3,2),(5,4)\)
\item
  \((\pi,\pi),(-\pi,-\pi)\)
\end{enumerate}

Model the following three problems using Cosine. Assume any point is either a maximum or a minimum point.

\begin{enumerate}
\def\labelenumi{\arabic{enumi}.}
\tightlist
\item
  \((2,2),(3,4)\)
\item
  \((0,1),(2\pi,3)\)
\item
  \((2\pi,2\pi),(-2\pi,-2\pi)\)
\end{enumerate}

\subsection{Topic 3: Graphing using Unit Circle}\label{topic-3-graphing-using-unit-circle}

\subsubsection{Graph the following problems}\label{graph-the-following-problems}

\begin{figure}[htbp]
	 \centering
	 \includegraphics[width=0.9\textwidth]{C:/Users/score/Downloads/graphingSinEZ.png}
\end{figure}

\begin{figure}[htbp]
	 \centering
	 \includegraphics[width=0.9\textwidth]{C:/Users/score/Downloads/graphingSinMED.png}
\end{figure}

\begin{figure}[htbp]
	 \centering
	 \includegraphics[width=0.9\textwidth]{C:/Users/score/Downloads/graphingSinHARD.png}
\end{figure}




\end{document}
